\chapter{Integration with Chainladder}
\label{chap:chainladder}

\section{Overview}

\actsim is designed to be seamlessly integrated with the established
triangle-based reserving tools. The \texttt{chainladder} Python package
provides a comprehensive set of methods for reserving, including
chain-ladder development, Bornhuetter--Ferguson, Cape Cod, Mack,
bootstrap chain-ladder, and various GLM-based approaches.

A natural division of labour is for \actsim to generate synthetic
claim-level data and convert the results into triangle format, while
\texttt{chainladder} performs reserving analysis on the resulting
triangle. This chapter shows how the two packages interact.

\section{From Simulated Claims to a Triangle}

Once \texttt{ClaimSimulator} has produced a long-format development
DataFrame, the data can be reshaped into the format expected by
\texttt{chainladder}.

\begin{lstlisting}[style=pythonstyle]
import pandas as pd
import numpy as np
from actsim import ClaimSimulator

# Build a policy table and simulate claims
policies = pd.DataFrame({
    "policy_id":  range(1, 501),
    "freq_dist":  "poisson",
    "freq_params": list(zip(np.random.uniform(0.6, 0.8, 500).round(2),)),
    "sev_dist":   "lognormal",
    "sev_params": list(zip(
        np.random.uniform(8, 12, 500).round(2),
        np.random.uniform(0.3, 0.7, 500).round(2),
    )),
    "start_date": pd.Timestamp("2018-01-01"),
    "end_date":   pd.Timestamp("2022-12-31"),
})

claim_sim = ClaimSimulator(policies, random_seed=42)
claim_sim.simulate_claims()
claim_sim.simulate_dates_nhpp(lambda0=10, alpha=0.5, phase=0, T=5)

base_LDFs = {0: 2.0, 3: 1.5, 6: 1.2, 9: 1.1, 12: 1.05, 15: 1.02, 18: 1.0}
claim_sim.simulate_claim_development(
    base_LDFs=base_LDFs,
    volatility=0.1,
    tail_factor=1.0,
)

triangle_df = (
    claim_sim.claim_development
    .groupby(["accident_year", "development_month"])["incurred_loss"]
    .sum()
    .reset_index()
)
\end{lstlisting}

The resulting \texttt{triangle\_df} has one row per
(accident year, development age) cell and aggregates the incurred losses
across all claims in that cell. Since the \texttt{chainladder}
package can also work directly with claim-level loss data and perform the
aggregation internally, this preprocessing step is optional.

\section{Constructing a chainladder Triangle}

The \texttt{chainladder} package accepts long-format data and converts
it into its native triangle representation:

\begin{lstlisting}[style=pythonstyle]
import chainladder as cl

triangle = cl.Triangle(
    triangle_df,
    origin="accident_year",
    development="development_month",
    columns="incurred_loss",
    cumulative=True,
)

print(triangle)
\end{lstlisting}

The \texttt{cumulative=True} flag indicates that the values represent
cumulative incurred losses; if instead the values represented
incremental amounts, the flag should be set to \texttt{False}.

\section{Position of ActSim Relative to Chainladder}

\actsim is positioned as a simulation and data-generation companion to
\texttt{chainladder} for analytical and scenario testing purposes. In a typical workflow:

\begin{itemize}[leftmargin=*]
    \item \actsim generates synthetic policy and claim-level data.
    \item \actsim develops the simulated claims through a user-supplied
          LDF schedule.
    \item The resulting development data are reshaped into a triangle.
    \item \texttt{chainladder} performs reserving analyses on the
          triangle.
    \item The reserving outputs are compared with the underlying
          simulated ultimates to validate the workflow.
\end{itemize}

This division of labour leverages the strengths of each package:
\actsim's flexible, transparent simulation engine and \texttt{chainladder}'s
comprehensive suite of reserving methods.
